% Filtration sous pression réduite sur entonnoir Büchner (L'essentiel, chapitre 9)
\begin{tikzpicture}[x=1cm,y=1cm]
% Fiole à vide et filtrat
\def\fiole{(-0.3,2.9) -- (-0.3,2.35) -- (-1.15,0.16) .. controls (-1.2,0.03) and (-1.12,0) .. (-1.0,0) -- (1.0,0)
  .. controls (1.12,0) and (1.2,0.03) .. (1.15,0.16) -- (0.3,2.35) -- (0.3,2.9)}
\begin{scope}\clip \fiole -- cycle; \fill[yellow!12] (-1.3,0) rectangle (1.3,0.6);\end{scope}
\draw[pcSurface] (-0.93,0.6) -- (0.93,0.6);
\draw[pcVerre] \fiole;
\draw[pcVerre] (0.3,2.62) -- (1.15,2.78) (0.3,2.48) -- (1.15,2.64);
\draw[otInk!70,line width=2.4pt,line cap=round] (1.15,2.71) -- (2.3,2.71);
% Bouchon et tige du Büchner
\filldraw[fill=otInk!65,draw=otInk!80,line width=0.4pt] (-0.36,2.9) -- (0.36,2.9) -- (0.3,3.18) -- (-0.3,3.18) -- cycle;
\draw[pcVerre] (-0.08,2.3) -- (-0.08,3.3) (0.08,2.3) -- (0.08,3.3);
% Entonnoir Büchner
\draw[pcVerre] (-0.08,3.3) -- (-0.75,3.55) -- (-0.75,4.45) (0.08,3.3) -- (0.75,3.55) -- (0.75,4.45);
\draw[otInk!55,line width=1.2pt] (-0.7,3.62) -- (0.7,3.62);
\foreach \x in {-0.55,-0.4,...,0.56}{\fill[otMuted!60] (\x,3.69) circle (0.05) (\x+0.07,3.78) circle (0.045);}
\pcNote[east]{(-0.75,4.2)}{(-1.4,4.4)}{entonnoir Büchner}
\pcNote[east]{(-0.6,3.68)}{(-1.4,3.65)}{cristaux sur papier filtre}
\pcNote[east]{(-0.33,3.04)}{(-1.4,2.9)}{bouchon}
\pcNote[east]{(-0.75,1.2)}{(-1.6,1.4)}{fiole à vide}
\pcNote[west]{(0.6,0.4)}{(1.6,0.55)}{filtrat}
\node[pcLegende,anchor=west] at (2.4,2.71) {vers la trompe à eau};
\end{tikzpicture}
\par\smallskip\legende{La trompe à eau aspire l'air de la fiole : la pression réduite accélère la filtration et essore les cristaux.}
